\documentclass{btu-thesis}
% Bu dosyayı kendi bilgilerine göre doldur.
% Başlık en fazla üç satır olacak şekilde \\ ile bölünebilir.

\renewcommand{\BTUThesisTitleTR}{BİTİRME ÇALIŞMASI BAŞLIĞI BURAYA GELİR\\GEREKLİ İSE İKİNCİ SATIR\\GEREKLİ İSE ÜÇÜNCÜ SATIR, ÜÇ SATIRA SIĞDIRINIZ}
\renewcommand{\BTUThesisTitleEN}{THESIS TITLE IN ENGLISH HERE}
\renewcommand{\BTUThesisTitleApproval}{BİTİRME ÇALIŞMASI BAŞLIĞI}

\renewcommand{\BTUStudentName}{Öğrenci Adı SOYADI}
\renewcommand{\BTUStudentNumber}{Öğrenci No}

\renewcommand{\BTUAdvisor}{Ünvan Ad SOYAD}
\renewcommand{\BTUAdvisorFull}{Prof. Dr. Adı SOYADI}
\renewcommand{\BTUJuryOne}{Dr. Öğr. Üyesi  Adı SOYADI}
\renewcommand{\BTUJuryTwo}{Öğr. Gör. Dr.  Adı SOYADI}
\renewcommand{\BTUChair}{Doç. Dr. Adı Soyadı}

\renewcommand{\BTUSubmissionMonthYear}{TEZİN SAVUNULDUĞU AY YIL}
\renewcommand{\BTUDefenseDate}{28 Haziran 2017}
\renewcommand{\BTUForewordDate}{Haziran 2017}

\renewcommand{\BTUKeywordsTR}{anahtar kelime 1, anahtar kelime 2, anahtar kelime 3}
\renewcommand{\BTUKeywordsEN}{keyword 1, keyword 2, keyword 3}

\usetikzlibrary{arrows.meta,positioning,calc,fit,backgrounds,shapes.geometric,decorations.pathreplacing}

\definecolor{BTUInk}{HTML}{1F2933}
\definecolor{BTUBlue}{HTML}{1F4E79}
\definecolor{BTUTeal}{HTML}{1B7F79}
\definecolor{BTUOrange}{HTML}{B7651B}
\definecolor{BTUGreen}{HTML}{4B7F2A}
\definecolor{BTUGray}{HTML}{6B7280}

\tikzset{
  thesisBox/.style={
    draw=BTUInk!55,
    line width=0.45pt,
    rounded corners=1.5pt,
    fill=BTUBlue!7,
    align=center,
    font=\scriptsize\sffamily,
    inner xsep=4pt,
    inner ysep=3pt,
    minimum height=8mm
  },
  thesisData/.style={thesisBox, fill=BTUTeal!9, draw=BTUTeal!70},
  thesisFeature/.style={thesisBox, fill=BTUOrange!10, draw=BTUOrange!70},
  thesisModel/.style={thesisBox, fill=BTUBlue!9, draw=BTUBlue!70},
  thesisEval/.style={thesisBox, fill=BTUGreen!10, draw=BTUGreen!70},
  thesisNote/.style={
    draw=BTUGray!45,
    line width=0.35pt,
    rounded corners=1.5pt,
    fill=BTUGray!6,
    align=center,
    font=\tiny\sffamily,
    inner xsep=3pt,
    inner ysep=2pt
  },
  thesisArrow/.style={-{Latex[length=2.2mm,width=1.5mm]}, line width=0.55pt, draw=BTUInk!75},
  thesisSoftArrow/.style={-{Latex[length=2mm,width=1.3mm]}, line width=0.45pt, draw=BTUGray!75},
  thesisDashedArrow/.style={thesisSoftArrow, dashed},
  graphNode/.style={circle, draw=BTUBlue!75, fill=BTUBlue!9, line width=0.5pt, minimum size=8mm, align=center, font=\tiny\sffamily},
  targetNode/.style={circle, draw=BTUGray!75, fill=BTUGray!10, line width=0.5pt, minimum size=8mm, align=center, font=\tiny\sffamily},
  obsNode/.style={draw=BTUTeal!75, fill=BTUTeal!9, line width=0.5pt, rounded corners=1.5pt, minimum width=17mm, minimum height=8mm, align=center, font=\tiny\sffamily},
  evidenceStep/.style={thesisBox, minimum width=28mm, minimum height=12mm}
}

\begin{document}

\BTUfrontmatter
\BTUmakecover
% Word şablonunda dış kapaktan sonra boş sayfa yorumu var. Bölüm özellikle isterse şu satırı aç:
% \thispagestyle{btunofooter}\null\clearpage
\BTUmakeinnercover
\BTUmakeapprovalpage
\BTUmakeplagiarism

\begin{BTUforeword}
Bu bitirme çalışmasının hazırlanma sürecinde bilgi birikimi, yönlendirmeleri ve yapıcı geri bildirimleriyle çalışmamıza katkı sunan danışmanımız Doç. Dr. Mustafa Özgür CİNGİZ'e teşekkür ederiz. Tez konusunun şekillenmesinden sonuçların değerlendirilmesine kadar gösterdiği akademik rehberlik, bu çalışmanın daha tutarlı ve savunulabilir bir çerçeveye kavuşmasını sağlamıştır.

Bursa Teknik Üniversitesi Bilgisayar Mühendisliği Bölümünde öğrenimimiz boyunca değerli bilgi ve deneyimlerini bizlerle paylaşan tüm hocalarımıza teşekkür ederiz. Lisans eğitimimiz süresince kazandığımız teknik bakış açısı, problem çözme disiplini ve araştırma kültürü bu çalışmanın temelini oluşturmuştur.

Bu süreçte sabırları, anlayışları ve sürekli destekleriyle yanımızda olan ailelerimize ve arkadaşlarımıza içten teşekkür ederiz. Yoğun çalışma dönemlerinde verdikleri moral ve güven, bu çalışmanın tamamlanmasında bizim için önemli bir güç olmuştur.

\BTUforewordsignature
\end{BTUforeword}

\BTUaddcontents

\begin{BTUabbr}{KISALTMALAR}
AUPRC & Precision-recall eğrisi alanı (area under precision-recall curve)\\
AUROC & ROC eğrisi alanı (area under receiver operating characteristic curve)\\
BCE & İkili çapraz entropi (binary cross-entropy)\\
BiLSTM & Çift yönlü LSTM (bidirectional long short-term memory)\\
CNN & Evrişimli sinir ağı (convolutional neural network)\\
CRISPR & Clustered regularly interspaced short palindromic repeats\\
CTCF & CCCTC bağlanma faktörü (CCCTC-binding factor)\\
DNA & Deoksiribonükleik asit (deoxyribonucleic acid)\\
DNase & Deoksiribonükleaz (deoxyribonuclease)\\
DRIP & DNA-RNA immünopresipitasyonu (DNA-RNA immunoprecipitation)\\
F1 & Kesinlik ve duyarlılığın harmonik ortalaması (F1 skoru)\\
FiLM & Özellik düzeyinde doğrusal modülasyon (feature-wise linear modulation)\\
FN & Yanlış negatif (false negative)\\
FP & Yanlış pozitif (false positive)\\
GAT & Çizge dikkat ağı (graph attention network)\\
GATv2 & Çizge dikkat ağı sürüm 2 (graph attention network version 2)\\
GCN & Çizge evrişimli ağ (graph convolutional network)\\
GNN & Çizge sinir ağı (graph neural network)\\
H3K4me3 & Histon H3 lizin 4 trimetilasyonu\\
MCC & Matthews korelasyon katsayısı\\
MNase & Mikrokokal nükleaz (micrococcal nuclease)\\
RNA & Ribonükleik asit (ribonucleic acid)\\
RNN & Tekrarlayan sinir ağı (recurrent neural network)\\
ROC & Alıcı işletim karakteristiği (receiver operating characteristic)\\
RRBS & İndirgenmiş temsilli bisülfit dizileme\\
SE & Sıkıştırma ve uyarım (squeeze-and-excitation)\\
sgRNA & Tek rehber RNA (single-guide RNA)\\
SHAP & Shapley katkılı açıklamalar (SHapley Additive exPlanations)\\
TN & Doğru negatif (true negative)\\
TP & Doğru pozitif (true positive)\\
XGBoost & Extreme gradient boosting\\
\end{BTUabbr}

\begin{BTUabbr}{SEMBOLLER}
\(\mathcal{D}_0\) & Ham aday veri evreni\\
\(\mathcal{D}_m\) & Deneysel olarak ölçülmüş veri evreni\\
\(\mathcal{G}_{tr}, \mathcal{G}_{te}\) & Eğitim ve test rehber kümeleri\\
\(\mathcal{G}_{A/B/C}\) & Graph A, Graph B ve Graph C çizge temsilleri\\
\(y\) & İkili hedef dışı aktivite etiketi\\
\(f_c\) & Ham kesim frekansı\\
\(X^{F1:F4}\) & Tabular feature ladder temsili\\
\(S_{g,t}\) & Rehber-hedef dizi çifti temsili\\
\(g_i\) & \(i\). rehber RNA düğümü\\
\(t_j\) & \(j\). fiziksel hedef düğümü\\
\(o_{ij}\) & \(g_i\) rehberiyle ilişkili hedef-gözlem düğümü\\
\(e_{ij}\) & \(g_i\) ve hedef/gözlem düğümü arasındaki aday kenar\\
\(x_{ij}^{edge}\) & Aday kenara ait özellik vektörü\\
\(x_{ij}^{obs}\) & Hedef-gözlem bağlam vektörü\\
\(x^{seq}, x^{epi}\) & Dizi ve epigenetik alt-aile vektörleri\\
\(x^{nuc}, x^{miss}\) & Nükleozom ve eksiklik alt-aile vektörleri\\
\(\phi_{\cdot}\) & Özellik alt-ailesi encoder fonksiyonu\\
\(z\) & Birleştirilmiş target-context temsili\\
\(\Gamma\) & Gating veya feature recalibration vektörü\\
\(h_g, h_o\) & Rehber ve hedef-gözlem gizil durumları\\
\(\alpha_{ij}^{(\ell)}\) & \(\ell\). katmanda \(i,j\) çifti için attention katsayısı\\
\(\hat{p}_{ij}\) & \(i,j\) adayı için tahmin edilen pozitif sınıf olasılığı\\
\(\oplus\) & Vektör birleştirme işlemi\\
\(\odot\) & Eleman bazlı çarpım\\
\(\mathrm{TP}, \mathrm{TN}\) & Doğru pozitif ve doğru negatif sayımları\\
\(\mathrm{FP}, \mathrm{FN}\) & Yanlış pozitif ve yanlış negatif sayımları\\
\end{BTUabbr}

\BTUaddtablelist
\BTUaddfigurelist

\begin{BTUabstractTR}
Türkçe özet 300 ile 500 kelime arasında olmalıdır. Özetlerde tezde ele alınan konu kısaca tanıtılır, kullanılan yöntemler ve ulaşılan sonuçlar belirtilir. Özetlerde kaynak, şekil ve çizelge verilmez. Bu paragraf şablonun derlenebilmesi için geçici olarak bırakılmıştır; gerçek tez metni yazılırken bu alan kendi özetinle değiştirilmelidir.
\end{BTUabstractTR}

\begin{BTUabstractEN}
The English summary must be between 300 and 500 words. It should briefly introduce the subject of the thesis, describe the methods used, and summarize the conclusions. References, figures, and tables must not be included in the summary. This paragraph is only a placeholder for compilation and should be replaced with the actual thesis summary.
\end{BTUabstractEN}

\BTUmainmatter
\chapter{GİRİŞ}
\label{chap:giris}

CRISPR-Cas9 sistemi, programlanabilir RNA rehberliğiyle DNA üzerinde belirli bölgelerin hedeflenebilmesini sağlayan temel bir genom düzenleme yaklaşımıdır. Cas9 enziminin RNA kılavuzlu hedef tanıma mekanizması, genom düzenleme çalışmalarında hedef bölge seçimini hesaplanabilir bir probleme dönüştürmüştür (Jinek ve diğ., 2012; Cong ve diğ., 2013). Bununla birlikte hedef sekansa benzeyen farklı genomik bölgelerde istenmeyen kesimlerin oluşabilmesi, CRISPR-Cas9 uygulamalarında hedef dışı bölge (off-target site) tahminini merkezi bir güvenilirlik sorunu haline getirmiştir (Hsu ve diğ., 2013).

Rehber RNA ve aday hedef DNA dizisi arasındaki eşleşme kalitesi, hedef dışı tahmin probleminin temel girdilerinden biridir. Aynı veya benzer sekans örüntüleri farklı kromatin bağlamlarında farklı erişilebilirlik ve aktivite davranışı gösterebilir. Cas9 bağlanması ve kesim etkinliğinin kromatin erişilebilirliği, nükleozom konumu ve hücresel bağlamla ilişkili olabileceği deneysel ve derleme çalışmalarında gösterilmiştir (Wu ve diğ., 2014; Horlbeck ve diğ., 2016; Isaac ve diğ., 2016; Yarrington ve diğ., 2018; Verkuijl ve Rots, 2019). Bu nedenle aday hedef bölgenin epigenetik durumu, nükleozom ilişkili özellikleri, deneysel ölçüm koşulu ve hedef gözleminin bağlamı tahmin probleminde ayrı bir bilgi ekseni oluşturur. crisprSQL veri tabanı CRISPR/Cas hedef dışı kesim deneylerini düzenli bir kaynak altında toplamış, Mak ve diğ. (2022) ise bu veri çizgisi üzerinde deneysel epigenetik ve hesaplanmış nükleozom tanımlayıcılarının hedef dışı aktiviteyle ilişkisini incelemiştir (Störtz ve Minary, 2021; Mak ve diğ., 2022).

Bu tezde incelenen problem, CRISPR-Cas9 hedef dışı tahmininde epigenetik ve kromatin bağlamının çizge tabanlı modeller içinde nasıl temsil edilmesi gerektiğidir. Çizge sinir ağları (graph neural networks, GNN), düğümler ve kenarlar arasında bilgi aktarımı yaparak ilişkisel yapıların öğrenilmesine olanak verir. Kipf ve Welling (2017) tarafından yaygınlaştırılan graph convolutional network (GCN) yaklaşımı, bu tür ilişkisel temsiller için temel bir model ailesi oluşturmuştur. CRISPR hedef dışı tahmininde sgRNA--hedef ilişkisini çizge bağlantısı olarak ele alan önceki çalışmalar, link prediction temelli GNN formülasyonunun bu problem alanına uygulanabilir olduğunu göstermiştir (Vinodkumar ve diğ., 2021). Bu tezde ise sgRNA--target bağlantı çizgesinin ötesine geçilerek, hedef gözlemine ait epigenetik ve kromatin bağlamının model içinde hangi temsille daha anlamlı kullanılabildiği sınanmıştır.

Çalışmanın değerlendirme ekseni ikiye ayrılmıştır. İlk eksen, eşikten bağımsız sıralama başarımıdır. Bu eksende area under the precision--recall curve (AUPRC) birincil metrik olarak kullanılmıştır. Dengesiz sınıf dağılımlarında precision--recall eğrileri, ROC eğrilerine göre pozitif sınıf başarımını daha doğrudan yansıtır (Davis ve Goadrich, 2006; Saito ve Rehmsmeier, 2015). İkinci eksen, validation üzerinde seçilmiş karar eşiğinde nadir ölçülmüş negatiflerin ayırt edilebilmesidir. Bu eksende Matthews correlation coefficient (MCC), macro-F1, specificity ve karışıklık matrisi değerleri ana yorum kanıtları olarak kullanılmıştır. Bu iki eksen aynı sonucu ölçmez. AUPRC modelin genel sıralama kalitesini değerlendirirken, karar eşiğine bağlı metrikler belirli bir kullanım noktasındaki sınıf ayırma davranışını gösterir.

Bu ayrım, çalışmanın ana iddia sınırını belirler. Bu tezde geliştirilen context-aware Graph C / GATv2 çizgisi, XGBoost F4 tabular referansına karşı sağlam bir AUPRC üstünlüğü olarak sunulmamaktadır. Bu iddia sınırı, Bölüm~\ref{chap:deneysel}'de verilen sağlamlık sonuçları ve paired-difference dağılımlarıyla desteklenmiştir (Çizelge~\ref{tab:auprc-saglamlik}; Şekil~\ref{fig:saglamlik-auprc-araliklari}; Şekil~\ref{fig:paired-difference-distributions}). Buna karşılık Graph C target-observation temsili, family-aware target-context encoder ve sequence/context fusion varyantları, validation-kilitli karar eşiğinde nadir ölçülmüş negatiflerin tanınması tarafında savunulabilir bir katkı göstermiştir. Bu katkının deneysel karşılığı Graph C GATv2, mekanizma çıkarma ve family-aware/context etkileşim sonuçlarında verilmiştir (Çizelge~\ref{tab:attention-sonuclari}; Çizelge~\ref{tab:graphc-mekanizma}; Çizelge~\ref{tab:family-aware-context}). Bu nedenle çalışma, en güçlü genel tahmin edici iddiası yerine, bağlam duyarlı çizge temsilinin nerede işe yaradığını ve nerede sınırlı kaldığını gösteren kontrollü bir modelleme ve değerlendirme çalışması olarak konumlandırılmıştır.

\section{Tezin Amacı}

Bu tezin amacı, CRISPR-Cas9 hedef dışı tahmininde epigenetik ve kromatin bağlamının çizge tabanlı modeller içinde temsil edilme biçimini incelemek ve bu temsilin hem sıralama başarımı hem de nadir negatif sınıf tanıma davranışı üzerindeki etkisini değerlendirmektir. Çalışmada Mak ve diğ. (2022) veri çizgisinden gelen crisprSQL türevi epigenetik-nükleozom veri seti kullanılmıştır. Ana hedef, aynı değerlendirme sözleşmesi altında tabular baseline, sequence tabanlı baseline, GCN, GAT/GATv2 ve Graph C target-observation temsillerinin karşılaştırılmasıdır.

Değerlendirme sözleşmesi, sonuçların güvenilir yorumlanabilmesi için dondurulmuştur. Ana ikili etiket Scheme A olarak tanımlanmış ve \texttt{cleavage\_freq > 1e-5} kuralı kullanılmıştır. NaN \texttt{cleavage\_freq} satırları supervised label üretiminden çıkarılmıştır. Ana veri ayrımı guide-disjoint olacak biçimde kurulmuştur. Ana validation ve test evreni measured-only satırlardan oluşmuş, \texttt{experiment\_id = 18} ana değerlendirme dışında bırakılmıştır. Ön işleme istatistikleri eğitim verisinden öğrenilmiş, model seçimi validation AUPRC ile yapılmış ve test seti üzerinde threshold veya model ayarı yapılmamıştır.

Bu çalışma kapsamında üç temsil düzeyi karşılaştırılmıştır. Graph A, fiziksel hedef düğümlerine dayanan minimal sgRNA--target çizgesi olarak kullanılmıştır. Graph B, rehber benzerliğiyle zenginleştirilmiş sınırlı bir kontrol olarak ele alınmıştır. Graph C ise her satırı ayrı bir target-observation düğümü olarak temsil ederek, aynı fiziksel hedef koordinatının farklı deneysel veya epigenetik bağlamlarda farklı gözlemler taşıyabilmesini model içine yansıtmıştır. Bu ayrım, Graph C sonucunun salt topology enrichment olarak yorumlanmaması için özellikle korunmuştur.

Tezin pratik amacı, güçlü bir tabular referans olan XGBoost F4 karşısında GNN modellerin koşulsuz üstün olduğunu göstermek değildir. Daha dar ve savunulabilir amaç, bağlam duyarlı Graph C GATv2 çizgisinin hangi koşullarda yararlı sinyal ürettiğini, hangi metriklerde sınırlı kaldığını ve bu etkinin hangi feature ailelerine bağlı olduğunu göstermektir. Bu nedenle AUPRC birincil metrik olarak korunmuş, MCC, macro-F1, specificity ve TN/FP/FN/TP değerleri ise nadir negatif operating-point davranışının ana kanıtları olarak değerlendirilmiştir.

\section{Literatür Araştırması}

CRISPR-Cas9 hedef dışı tahmin literatüründe ilk ana eksen, rehber RNA ve hedef DNA sekansları arasındaki benzerliğe dayalı modellemedir. Hsu ve diğ. (2013), Cas9 hedefleme özgüllüğünün rehber--hedef uyumsuzluk örüntülerinden güçlü biçimde etkilendiğini göstermiştir. Daha sonraki derin öğrenme yaklaşımları, sekans çiftlerini CNN, RNN, attention, çoklu görünüm veya transformer tabanlı temsilcilerle işleyerek hedef dışı aktiviteyi tahmin etmeye çalışmıştır (Chuai ve diğ., 2018; Lin ve diğ., 2020; Liu ve diğ., 2020; Luo ve diğ., 2024; Sun ve diğ., 2024). Bu çalışmalar, sekans temsilinin problem için zorunlu bir temel olduğunu göstermektedir. Ancak farklı veri kaynakları, negatif örnekleme stratejileri, guide split politikaları ve sınıf prevalansları nedeniyle bu literatürdeki ham skorlar bu tezdeki measured-only test evreniyle doğrudan karşılaştırılamaz.

İkinci eksen, epigenetik ve kromatin bağlamının hedef dışı aktiviteyle ilişkisidir. crisprSQL, farklı CRISPR/Cas hedef dışı kesim deneylerini yapılandırılmış biçimde bir araya getirmiştir (Störtz ve Minary, 2021). Mak ve diğ. (2022), bu veri çizgisi üzerinde 6 deneysel epigenetik özellik ve 13 hesaplanmış nükleozom ilişkili özelliği incelemiştir. piCRISPR çalışması da aynı veri ailesine yakın bir hatta fiziksel olarak bilgilendirilmiş sequence, kromatin ve nükleozom tanımlayıcılarının hedef dışı kesim tahminindeki rolünü incelemiştir (Störtz ve diğ., 2023). Bu tez, aynı veri ve feature lineage'ından yararlansa da Mak ve diğ. (2022) çalışmasının birebir reprodüksiyonu değildir. Mak çalışması continuous CA dönüşümü, korelasyon ve SHAP temelli feature analizi etrafında yapılandırılmıştır. Bu tezde ise Scheme A ikili sınıflandırma hedefi, guide-disjoint split, measured-only validation/test evreni ve AUPRC birincil metriği kullanılmıştır.

Üçüncü eksen, CRISPR hedef dışı tahmininin çizge problemi olarak ele alınmasıdır. GCN mimarisi, komşuluk ilişkilerinden düğüm temsili öğrenmek için temel bir GNN yaklaşımı sağlamıştır (Kipf ve Welling, 2017). Vinodkumar ve diğ. (2021), sgRNA hedef dışı aktivitesinin GCN tabanlı bir yaklaşımla modellenebileceğini göstermiştir. Bu tezde bu çizgi, daha katı bir measured-only ve guide-disjoint değerlendirme sözleşmesiyle genişletilmiştir. Minimal fiziksel hedef çizgesi, guide-similarity kontrolü ve target-observation context çizgesi ayrı ayrı test edilmiştir. Böylece graph schema seçiminin ve hedef bağlamı semantiğinin model davranışı üzerindeki etkisi ayrıştırılmıştır.

Dördüncü eksen, attention tabanlı GNN mimarileridir. GAT modelleri, komşu düğümlerden gelen mesajları attention katsayılarıyla ağırlıklandırır (Veličković ve diğ., 2018). GATv2, attention mekanizmasının daha dinamik bir biçimini hedeflediği için Graph C bağlamında anlamlı bir model adayıdır (Brody ve diğ., 2022). Ancak bu tezde attention mekanizması tek başına başarı garantisi olarak ele alınmamıştır. Attention tabanlı modeller Graph A/B/C şemaları ve edge-aware kullanım koşulları altında kontrollü biçimde sınanmıştır. Graph C target-observation bağlamı ve candidate-edge enerji özelliklerinin birlikte kullanılması, nadir negatif operating-point davranışını güçlendiren ana yön olarak değerlendirilmiştir.

Beşinci eksen, değerlendirme metodolojisi ve sınıf dengesizliği yorumudur. Precision--recall eğrileri, sınıf dağılımının dengesiz olduğu durumlarda model karşılaştırmasını ROC eğrilerine göre daha anlamlı hale getirebilir (Davis ve Goadrich, 2006; Saito ve Rehmsmeier, 2015). Precision değerinin sınıf prevalansına duyarlı olması, AUPRC değerlerinin test evrenindeki sınıf oranıyla birlikte yorumlanmasını gerektirir (Williams, 2021). Bu tezde test evreni pozitif ağırlıklıdır ve pozitif prevalans 0.900705'tir. Bu nedenle problem, genom çapında negatif ağırlıklı retrieval benchmark'larından farklıdır. Nadir sınıf bu ölçüm evreninde negatiftir. Bu durum, specificity, MCC, macro-F1 ve karışıklık matrisinin yorum değerini artırmıştır. Sağlamlık analizinde guide-cluster bootstrap, paired-difference analizleri ve çoklu tohum değerlendirmesi kullanılarak tekil eğitim koşulundaki AUPRC kazanımlarının belirsizlik içinde kalıp kalmadığı ayrıca incelenmiştir (Field ve Welsh, 2007; Bethard, 2022).

Bu literatür içinde tezin konumu, sekans odaklı bir karşılaştırma çalışması veya Mak ve diğ. (2022) reprodüksiyonu değildir. Çalışma, epigenetik/kromatin bağlamının çizge temsiliyle nasıl kullanılabileceğine odaklanır. Bu nedenle literatürden alınan fikirler birebir yeniden üretim iddiasıyla değil, yerel veri sözleşmesi altında uyarlanmış modelleme kararları olarak ele alınmıştır.

\section{Hipotez}

Bu tezde üç temel hipotez sınanmıştır.

Birinci hipotez, Graph C target-observation context temsili ve family-aware GATv2 encoder'ın, sabit measured-only ve guide-disjoint değerlendirme sözleşmesi altında nadir negatif sınıf tanıma davranışını Graph A GCN ve düz Graph C GCN referanslarına göre iyileştireceğidir. Bu hipotez, AUPRC yerine validation-kilitli threshold metrikleriyle ilişkilidir. MCC, macro-F1, specificity ve TN/FP/FN/TP değerleri bu davranışın ana kanıtları olarak değerlendirilmiştir. Hipotezin deneysel karşılığı Graph C GATv2 ve family-aware encoder sonuçlarında izlenmiştir (Çizelge~\ref{tab:attention-sonuclari}; Çizelge~\ref{tab:family-aware-context}; Şekil~\ref{fig:sequence-context-threshold}).

İkinci hipotez, context-aware GNN çizgisinin XGBoost F4'e karşı sağlam bir AUPRC üstünlüğü oluşturmayacağıdır. Bu hipotez, çalışmanın iddia sınırını açık biçimde belirler. XGBoost F4, ilk dondurulmuş değerlendirmede test AUPRC 0.992522 ile en güçlü tabular referans olarak kalmıştır (Çizelge~\ref{tab:non-graph-referanslar}). Sağlamlık analizinde XGBoost version drift nedeniyle yeniden üretilen F4 AUPRC değeri 0.992338 olarak raporlanmış, paired AUPRC farkları ise sıfırı dışlamamıştır (Çizelge~\ref{tab:auprc-saglamlik}; Şekil~\ref{fig:paired-difference-distributions}). Bu nedenle GNN sonuçları rekabetçi fakat AUPRC üstünlüğü iddiası taşımayan sonuçlar olarak yorumlanmıştır.

Üçüncü hipotez, Graph C kazanımının explicit context-similarity edge'lerinden çok doğrudan target-observation context feature'larına, özellikle experimental epigenetic feature ailesine bağlı olduğudur. Mekanizma ve alt-özellik maskeleme analizleri bu hipotezin temel deneysel karşılığıdır. Bu kanıtlar Çizelge~\ref{tab:graphc-mekanizma}, Çizelge~\ref{tab:target-context-mask} ve Şekil~\ref{fig:target-context-subgroup-threshold}'de verilmiştir. Target-observation context feature'ları maskelendiğinde model davranışının belirgin biçimde bozulması, Graph C katkısının ek topolojiyle tek başına açıklanamayacağını göstermiştir. Bu bulgu biyolojik olarak kromatin ve epigenetik bağlam literatürüyle uyumludur; ancak attention ağırlıkları, feature masking sonuçları veya context gate davranışları nedensel biyolojik kanıt olarak yorumlanmamıştır.

Bu hipotezler birlikte ele alındığında tez, CRISPR-Cas9 hedef dışı tahmininde context-aware GNN yaklaşımının değerini birincil AUPRC üstünlüğüyle değil, kontrollü değerlendirme altında mekanizma izole edilmiş ve validation-kilitli nadir negatif tanıma davranışıyla ortaya koymayı amaçlamaktadır.

\chapter{MATERYAL VE YÖNTEM}

Şekil ve çizelgeler metinde ilk kez anıldıktan sonra uygun en yakın yere yerleştirilmelidir. Örnek olarak Şekil~\ref{fig:ornek-sekil} ve Çizelge~\ref{tab:ornek-cizelge} verilmiştir.

\begin{figure}[htbp]
\centering
\fbox{\rule{0pt}{3cm}\rule{6cm}{0pt}}
\caption{Tek satır olması durumunda ortalı olacak.}
\label{fig:ornek-sekil}
\end{figure}

\begin{table}[htbp]
\caption{Tek satırlı ve kolonlar ortalanmış çizelge.}
\label{tab:ornek-cizelge}
\centering
\begin{tabular}{c c c c}
\toprule\toprule
Kolon A & Kolon B & Kolon C & Kolon D\\
\midrule
Satır A & Satır A & Satır A & Satır A\\
Satır B & Satır B & Satır B & Satır B\\
Satır C & Satır C & Satır C & Satır C\\
\bottomrule
\end{tabular}
\end{table}

\section{Denklemler}
Denklemler 12 punto yazılır ve denklem numaraları sağa dayalı olur. Örneğin Denklem~\ref{eq:ornek} aşağıda verilmiştir.

\begin{equation}
E = mc^2
\label{eq:ornek}
\end{equation}

Parametreler denklemin altında açıklanmalıdır.

\chapter{DENEYSEL ÇALIŞMALAR}

Bu bölüm deneysel çalışmalar, uygulama, modelleme veya analiz bölümü olarak yeniden adlandırılabilir.

\section{Çalışmanın Uygulama Alanı}
Bu kısımda uygulama alanı açıklanır.

\chapter{SONUÇ VE ÖNERİLER}

Bu bölümde çalışmanın sonuçları ve öneriler verilir.


\begin{BTUreferences}\sloppy
\item Bengio, Y. \& Grandvalet, Y. (2004). No unbiased estimator of the variance of k-fold cross-validation. \textit{Journal of Machine Learning Research, 5}, 1089--1105.
\item Bethard, S. (2022). We need to talk about random seeds. \textit{arXiv}. \url{https://arxiv.org/abs/2210.13393}
\item Boyd, K., Eng, K. H. \& Page, C. D. (2013). Area under the precision-recall curve: Point estimates and confidence intervals. In H. Blockeel, K. Kersting, S. Nijssen \& F. Zelezny (Eds.), \textit{Machine learning and knowledge discovery in databases} (Lecture Notes in Computer Science, Vol. 8190, pp. 451--466). Springer. \url{https://doi.org/10.1007/978-3-642-40994-3_29}
\item Brody, S., Alon, U. \& Yahav, E. (2022). How attentive are graph attention networks? \textit{International Conference on Learning Representations}. \url{https://openreview.net/forum?id=F72ximsx7C1}
\item Brockschmidt, M. (2020). GNN-FiLM: Graph neural networks with feature-wise linear modulation. In H. Daumé III \& A. Singh (Eds.), \textit{Proceedings of the 37th International Conference on Machine Learning} (Proceedings of Machine Learning Research, Vol. 119, pp. 1144--1152). PMLR. \url{https://proceedings.mlr.press/v119/brockschmidt20a.html}
\item Chen, T. \& Guestrin, C. (2016). XGBoost: A scalable tree boosting system. In \textit{Proceedings of the 22nd ACM SIGKDD International Conference on Knowledge Discovery and Data Mining} (pp. 785--794). Association for Computing Machinery. \url{https://doi.org/10.1145/2939672.2939785}
\item Chuai, G., Ma, H., Yan, J., Chen, M., Hong, N., Xue, D., Zhou, C., Zhu, C., Chen, K., Duan, B., Gu, F., Qu, S., Huang, D., Wei, J. \& Liu, Q. (2018). DeepCRISPR: Optimized CRISPR guide RNA design by deep learning. \textit{Genome Biology, 19}(1), 80. \url{https://doi.org/10.1186/s13059-018-1459-4}
\item Cong, L., Ran, F. A., Cox, D., Lin, S., Barretto, R., Habib, N., Hsu, P. D., Wu, X., Jiang, W., Marraffini, L. A. \& Zhang, F. (2013). Multiplex genome engineering using CRISPR/Cas systems. \textit{Science, 339}(6121), 819--823. \url{https://doi.org/10.1126/science.1231143}
\item Cui, Y., Jia, M., Lin, T.-Y., Song, Y. \& Belongie, S. (2019). Class-balanced loss based on effective number of samples. In \textit{Proceedings of the IEEE/CVF Conference on Computer Vision and Pattern Recognition} (pp. 9260--9269). \url{https://doi.org/10.1109/CVPR.2019.00949}
\item Davis, J. \& Goadrich, M. (2006). The relationship between precision-recall and ROC curves. In \textit{Proceedings of the 23rd International Conference on Machine Learning} (pp. 233--240). \url{https://doi.org/10.1145/1143844.1143874}
\item Field, C. A. \& Welsh, A. H. (2007). Bootstrapping clustered data. \textit{Journal of the Royal Statistical Society: Series B (Statistical Methodology), 69}(3), 369--390. \url{https://doi.org/10.1111/j.1467-9868.2007.00593.x}
\item Gao, Y., Chuai, G., Yu, W., Qu, S. \& Liu, Q. (2020). Data imbalance in CRISPR off-target prediction. \textit{Briefings in Bioinformatics, 21}(4), 1448--1454. \url{https://doi.org/10.1093/bib/bbz069}
\item Guan, Z. \& Jiang, Z. (2024). A systematic method for solving data imbalance in CRISPR off-target prediction tasks. \textit{Computers in Biology and Medicine, 178}, 108781. \url{https://doi.org/10.1016/j.compbiomed.2024.108781}
\item Horlbeck, M. A., Witkowsky, L. B., Guglielmi, B., Replogle, J. M., Gilbert, L. A., Villalta, J. E., Torigoe, S. E., Tjian, R. \& Weissman, J. S. (2016). Nucleosomes impede Cas9 access to DNA in vivo and in vitro. \textit{eLife, 5}, e12677. \url{https://doi.org/10.7554/eLife.12677}
\item Hsu, P. D., Scott, D. A., Weinstein, J. A., Ran, F. A., Konermann, S., Agarwala, V., Li, Y., Fine, E. J., Wu, X., Shalem, O. \& Zhang, F. (2013). DNA targeting specificity of RNA-guided Cas9 nucleases. \textit{Nature Biotechnology, 31}(9), 827--832. \url{https://doi.org/10.1038/nbt.2647}
\item Hu, J., Shen, L., Albanie, S., Sun, G. \& Wu, E. (2018). Squeeze-and-excitation networks. In \textit{Proceedings of the IEEE/CVF Conference on Computer Vision and Pattern Recognition} (pp. 7132--7141). \url{https://doi.org/10.1109/CVPR.2018.00745}
\item Isaac, R. S., Jiang, F., Doudna, J. A., Lim, W. A., Narlikar, G. J. \& Almeida, R. (2016). Nucleosome breathing and remodeling constrain CRISPR-Cas9 function. \textit{eLife, 5}, e13450. \url{https://doi.org/10.7554/eLife.13450}
\item Jain, S. \& Wallace, B. C. (2019). Attention is not explanation. In \textit{Proceedings of the 2019 Conference of the North American Chapter of the Association for Computational Linguistics: Human Language Technologies, Volume 1} (pp. 3543--3556). Association for Computational Linguistics. \url{https://doi.org/10.18653/v1/N19-1357}
\item Jinek, M., Chylinski, K., Fonfara, I., Hauer, M., Doudna, J. A. \& Charpentier, E. (2012). A programmable dual-RNA-guided DNA endonuclease in adaptive bacterial immunity. \textit{Science, 337}(6096), 816--821. \url{https://doi.org/10.1126/science.1225829}
\item Kapoor, S. \& Narayanan, A. (2023). Leakage and the reproducibility crisis in machine-learning-based science. \textit{Patterns, 4}(9), 100804. \url{https://doi.org/10.1016/j.patter.2023.100804}
\item Kipf, T. N. \& Welling, M. (2017). Semi-supervised classification with graph convolutional networks. \textit{International Conference on Learning Representations}. \url{https://openreview.net/forum?id=SJU4ayYgl}
\item Lin, T.-Y., Goyal, P., Girshick, R., He, K. \& Dollár, P. (2017). Focal loss for dense object detection. In \textit{Proceedings of the IEEE International Conference on Computer Vision} (pp. 2999--3007). \url{https://doi.org/10.1109/ICCV.2017.324}
\item Lin, J., Zhang, Z., Zhang, S., Chen, J. \& Wong, K.-C. (2020). CRISPR-Net: A recurrent convolutional network quantifies CRISPR off-target activities with mismatches and indels. \textit{Advanced Science, 7}(13), 1903562. \url{https://doi.org/10.1002/advs.201903562}
\item Liu, Q., Cheng, X., Liu, G., Li, B. \& Liu, X. (2020). Deep learning improves the ability of sgRNA off-target propensity prediction. \textit{BMC Bioinformatics, 21}, 51. \url{https://doi.org/10.1186/s12859-020-3395-z}
\item Luo, Y., Chen, Y., Xie, H., Zhu, W. \& Zhang, G. (2024). Interpretable CRISPR/Cas9 off-target activities with mismatches and indels prediction using BERT. \textit{Computers in Biology and Medicine, 169}, 107932. \url{https://doi.org/10.1016/j.compbiomed.2024.107932}
\item Mak, J. K., Störtz, F. \& Minary, P. (2022). Comprehensive computational analysis of epigenetic descriptors affecting CRISPR-Cas9 off-target activity. \textit{BMC Genomics, 23}, 805. \url{https://doi.org/10.1186/s12864-022-09012-7}
\item Perez, E., Strub, F., de Vries, H., Dumoulin, V. \& Courville, A. (2018). FiLM: Visual reasoning with a general conditioning layer. \textit{Proceedings of the AAAI Conference on Artificial Intelligence, 32}(1). \url{https://doi.org/10.1609/aaai.v32i1.11671}
\item Saito, T. \& Rehmsmeier, M. (2015). The precision-recall plot is more informative than the ROC plot when evaluating binary classifiers on imbalanced datasets. \textit{PLOS ONE, 10}(3), e0118432. \url{https://doi.org/10.1371/journal.pone.0118432}
\item Schenker, N. \& Gentleman, J. F. (2001). On judging the significance of differences by examining the overlap between confidence intervals. \textit{The American Statistician, 55}(3), 182--186. \url{https://doi.org/10.1198/000313001317097960}
\item Serrano, S. \& Smith, N. A. (2019). Is attention interpretable? In \textit{Proceedings of the 57th Annual Meeting of the Association for Computational Linguistics} (pp. 2931--2951). Association for Computational Linguistics. \url{https://doi.org/10.18653/v1/P19-1282}
\item Störtz, F., Mak, J. K. \& Minary, P. (2023). piCRISPR: Physically informed deep learning models for CRISPR/Cas9 off-target cleavage prediction. \textit{Artificial Intelligence in the Life Sciences, 3}, 100075. \url{https://doi.org/10.1016/j.ailsci.2023.100075}
\item Störtz, F. \& Minary, P. (2021). crisprSQL: A novel database platform for CRISPR/Cas off-target cleavage assays. \textit{Nucleic Acids Research, 49}(D1), D855--D861. \url{https://doi.org/10.1093/nar/gkaa885}
\item Sudre, C. H., Li, W., Vercauteren, T., Ourselin, S. \& Cardoso, M. J. (2017). Generalised Dice overlap as a deep learning loss function for highly unbalanced segmentations. In \textit{Deep learning in medical image analysis and multimodal learning for clinical decision support} (Lecture Notes in Computer Science, Vol. 10553, pp. 240--248). Springer. \url{https://doi.org/10.1007/978-3-319-67558-9_28}
\item Sun, J., Guo, J. \& Liu, J. (2024). CRISPR-M: Predicting sgRNA off-target effect using a multi-view deep learning network. \textit{PLOS Computational Biology, 20}(3), e1011972. \url{https://doi.org/10.1371/journal.pcbi.1011972}
\item Verkuijl, S. A. N. \& Rots, M. G. (2019). The influence of eukaryotic chromatin state on CRISPR-Cas9 editing efficiencies. \textit{Current Opinion in Biotechnology, 55}, 68--73. \url{https://doi.org/10.1016/j.copbio.2018.07.005}
\item Veličković, P., Cucurull, G., Casanova, A., Romero, A., Liò, P. \& Bengio, Y. (2018). Graph attention networks. \textit{International Conference on Learning Representations}. \url{https://openreview.net/forum?id=rJXMpikCZ}
\item Vinodkumar, P. K., Ozcinar, C. \& Anbarjafari, G. (2021). Prediction of sgRNA off-target activity in CRISPR/Cas9 gene editing using graph convolution network. \textit{Entropy, 23}(5), 608. \url{https://doi.org/10.3390/e23050608}
\item Williams, C. K. I. (2021). The effect of class imbalance on precision-recall curves. \textit{Neural Computation, 33}(4), 853--857. \url{https://doi.org/10.1162/neco_a_01362}
\item Wu, X., Scott, D. A., Kriz, A. J., Chiu, A. C., Hsu, P. D., Dadon, D. B., Cheng, A. W., Trevino, A. E., Konermann, S., Tsai, S. Q., Delebecque, N. J., Niemi, D. M., Long, C., Jaenisch, R., Zhang, F. \& Sharp, P. A. (2014). Genome-wide binding of the CRISPR endonuclease Cas9 in mammalian cells. \textit{Nature Biotechnology, 32}(7), 670--676. \url{https://doi.org/10.1038/nbt.2889}
\item Yarrington, R. M., Verma, S., Schwartz, S., Trautman, J. K. \& Carroll, D. (2018). Nucleosomes inhibit target cleavage by CRISPR-Cas9 in vivo. \textit{Proceedings of the National Academy of Sciences, 115}(38), 9351--9358. \url{https://doi.org/10.1073/pnas.1810062115}
\end{BTUreferences}

\BTUappendices
\refstepcounter{chapter}
\chapter*{EK A}
Ekler bölümünde verilen şekil ve çizelgeler, bulundukları ekin adı altında numaralandırılabilir. LaTeX varsayılanı bu bölümde A.1, A.2 şeklindedir.

\begin{figure}[htbp]
\centering
\fbox{\rule{0pt}{3cm}\rule{6cm}{0pt}}
\caption{Ekler bölümünde şekil örneği.}
\label{fig:ek-ornek}
\end{figure}

\begin{table}[htbp]
\caption{Ekler bölümünde çizelge örneği.}
\label{tab:ek-ornek}
\centering
\begin{tabular}{c c c c}
\toprule\toprule
Kolon A & Kolon B & Kolon C & Kolon D\\
\midrule
Satır A & Satır A & Satır A & Satır A\\
Satır B & Satır B & Satır B & Satır B\\
Satır C & Satır C & Satır C & Satır C\\
\bottomrule
\end{tabular}
\end{table}


\begin{BTUcv}
\begin{tabular}{@{}p{4.3cm}@{ : }p{9cm}@{}}
Ad-Soyad & \\
Doğum Tarihi ve Yeri & \\
E-posta & \\
\end{tabular}

\vspace{1cm}
\textbf{BİTİRME ÇALIŞMASINDAN TÜRETİLEN MAKALE, BİLDİRİ VEYA SUNUMLAR:}

\begin{itemize}
\item \ldots
\item \ldots
\end{itemize}
\end{BTUcv}

\end{document}
